% Options for packages loaded elsewhere
\PassOptionsToPackage{unicode}{hyperref}
\PassOptionsToPackage{hyphens}{url}
\documentclass[
]{article}
\usepackage{xcolor}
\usepackage[margin=1in]{geometry}
\usepackage{amsmath,amssymb}
\setcounter{secnumdepth}{-\maxdimen} % remove section numbering
\usepackage{iftex}
\ifPDFTeX
  \usepackage[T1]{fontenc}
  \usepackage[utf8]{inputenc}
  \usepackage{textcomp} % provide euro and other symbols
\else % if luatex or xetex
  \usepackage{unicode-math} % this also loads fontspec
  \defaultfontfeatures{Scale=MatchLowercase}
  \defaultfontfeatures[\rmfamily]{Ligatures=TeX,Scale=1}
\fi
\usepackage{lmodern}
\ifPDFTeX\else
  % xetex/luatex font selection
\fi
% Use upquote if available, for straight quotes in verbatim environments
\IfFileExists{upquote.sty}{\usepackage{upquote}}{}
\IfFileExists{microtype.sty}{% use microtype if available
  \usepackage[]{microtype}
  \UseMicrotypeSet[protrusion]{basicmath} % disable protrusion for tt fonts
}{}
\makeatletter
\@ifundefined{KOMAClassName}{% if non-KOMA class
  \IfFileExists{parskip.sty}{%
    \usepackage{parskip}
  }{% else
    \setlength{\parindent}{0pt}
    \setlength{\parskip}{6pt plus 2pt minus 1pt}}
}{% if KOMA class
  \KOMAoptions{parskip=half}}
\makeatother
\usepackage{color}
\usepackage{fancyvrb}
\newcommand{\VerbBar}{|}
\newcommand{\VERB}{\Verb[commandchars=\\\{\}]}
\DefineVerbatimEnvironment{Highlighting}{Verbatim}{commandchars=\\\{\}}
% Add ',fontsize=\small' for more characters per line
\usepackage{framed}
\definecolor{shadecolor}{RGB}{248,248,248}
\newenvironment{Shaded}{\begin{snugshade}}{\end{snugshade}}
\newcommand{\AlertTok}[1]{\textcolor[rgb]{0.94,0.16,0.16}{#1}}
\newcommand{\AnnotationTok}[1]{\textcolor[rgb]{0.56,0.35,0.01}{\textbf{\textit{#1}}}}
\newcommand{\AttributeTok}[1]{\textcolor[rgb]{0.13,0.29,0.53}{#1}}
\newcommand{\BaseNTok}[1]{\textcolor[rgb]{0.00,0.00,0.81}{#1}}
\newcommand{\BuiltInTok}[1]{#1}
\newcommand{\CharTok}[1]{\textcolor[rgb]{0.31,0.60,0.02}{#1}}
\newcommand{\CommentTok}[1]{\textcolor[rgb]{0.56,0.35,0.01}{\textit{#1}}}
\newcommand{\CommentVarTok}[1]{\textcolor[rgb]{0.56,0.35,0.01}{\textbf{\textit{#1}}}}
\newcommand{\ConstantTok}[1]{\textcolor[rgb]{0.56,0.35,0.01}{#1}}
\newcommand{\ControlFlowTok}[1]{\textcolor[rgb]{0.13,0.29,0.53}{\textbf{#1}}}
\newcommand{\DataTypeTok}[1]{\textcolor[rgb]{0.13,0.29,0.53}{#1}}
\newcommand{\DecValTok}[1]{\textcolor[rgb]{0.00,0.00,0.81}{#1}}
\newcommand{\DocumentationTok}[1]{\textcolor[rgb]{0.56,0.35,0.01}{\textbf{\textit{#1}}}}
\newcommand{\ErrorTok}[1]{\textcolor[rgb]{0.64,0.00,0.00}{\textbf{#1}}}
\newcommand{\ExtensionTok}[1]{#1}
\newcommand{\FloatTok}[1]{\textcolor[rgb]{0.00,0.00,0.81}{#1}}
\newcommand{\FunctionTok}[1]{\textcolor[rgb]{0.13,0.29,0.53}{\textbf{#1}}}
\newcommand{\ImportTok}[1]{#1}
\newcommand{\InformationTok}[1]{\textcolor[rgb]{0.56,0.35,0.01}{\textbf{\textit{#1}}}}
\newcommand{\KeywordTok}[1]{\textcolor[rgb]{0.13,0.29,0.53}{\textbf{#1}}}
\newcommand{\NormalTok}[1]{#1}
\newcommand{\OperatorTok}[1]{\textcolor[rgb]{0.81,0.36,0.00}{\textbf{#1}}}
\newcommand{\OtherTok}[1]{\textcolor[rgb]{0.56,0.35,0.01}{#1}}
\newcommand{\PreprocessorTok}[1]{\textcolor[rgb]{0.56,0.35,0.01}{\textit{#1}}}
\newcommand{\RegionMarkerTok}[1]{#1}
\newcommand{\SpecialCharTok}[1]{\textcolor[rgb]{0.81,0.36,0.00}{\textbf{#1}}}
\newcommand{\SpecialStringTok}[1]{\textcolor[rgb]{0.31,0.60,0.02}{#1}}
\newcommand{\StringTok}[1]{\textcolor[rgb]{0.31,0.60,0.02}{#1}}
\newcommand{\VariableTok}[1]{\textcolor[rgb]{0.00,0.00,0.00}{#1}}
\newcommand{\VerbatimStringTok}[1]{\textcolor[rgb]{0.31,0.60,0.02}{#1}}
\newcommand{\WarningTok}[1]{\textcolor[rgb]{0.56,0.35,0.01}{\textbf{\textit{#1}}}}
\usepackage{graphicx}
\makeatletter
\newsavebox\pandoc@box
\newcommand*\pandocbounded[1]{% scales image to fit in text height/width
  \sbox\pandoc@box{#1}%
  \Gscale@div\@tempa{\textheight}{\dimexpr\ht\pandoc@box+\dp\pandoc@box\relax}%
  \Gscale@div\@tempb{\linewidth}{\wd\pandoc@box}%
  \ifdim\@tempb\p@<\@tempa\p@\let\@tempa\@tempb\fi% select the smaller of both
  \ifdim\@tempa\p@<\p@\scalebox{\@tempa}{\usebox\pandoc@box}%
  \else\usebox{\pandoc@box}%
  \fi%
}
% Set default figure placement to htbp
\def\fps@figure{htbp}
\makeatother
\setlength{\emergencystretch}{3em} % prevent overfull lines
\providecommand{\tightlist}{%
  \setlength{\itemsep}{0pt}\setlength{\parskip}{0pt}}
\usepackage{bookmark}
\IfFileExists{xurl.sty}{\usepackage{xurl}}{} % add URL line breaks if available
\urlstyle{same}
\hypersetup{
  pdftitle={Regularisation: ridge{,} lasso{,} elastic net},
  hidelinks,
  pdfcreator={LaTeX via pandoc}}

\title{Regularisation: ridge, lasso, elastic net}
\author{}
\date{\vspace{-2.5em}}

\begin{document}
\maketitle

\textbf{Workflow labs} use the variant template: Goal → Approach →
Execution → Check → Report.

\subsection{Learning objectives}\label{learning-objectives}

\begin{itemize}
\tightlist
\item
  Fit ridge, lasso, and elastic-net regressions with \texttt{glmnet},
  and explain what each penalty does to the coefficient path.
\item
  Select \texttt{lambda} and \texttt{alpha} by cross-validation, and
  describe the bias--variance trade-off the choice implies.
\item
  Interpret a coefficient path figure and a set of standardised
  coefficients.
\end{itemize}

\subsection{Prerequisites}\label{prerequisites}

OLS and GLMs; the cross-validation lab.

\subsection{Background}\label{background}

The ordinary least-squares solution is unbiased but, when \emph{p} is
close to or larger than \emph{n}, it is unstable: small perturbations of
the data move the fitted coefficients a long way. Regularisation trades
a little bias for a lot of variance reduction by adding a penalty to the
fitting criterion. Ridge shrinks all coefficients toward zero, lasso
sets some exactly to zero, and elastic net blends the two. Each lives on
the \texttt{glmnet} coefficient path and is selected by the same basic
mechanism: cross-validation over \texttt{lambda} (and optionally
\texttt{alpha}).

In biomedical data the motivation is rarely pure prediction. A lasso
that selects 7 of 400 genes is also a shortlist for the next experiment.
A ridge that keeps every variable but shrinks their effects is an honest
admission that signal is diffuse and that no small subset will carry the
story. Elastic net is the default when variables come in correlated
groups, because the lasso tends to pick one at random from a correlated
cluster while ridge keeps them all.

Standardise before fitting. The penalty is scale-sensitive, and a
variable in raw units will be penalised differently from the same
variable in standardised units. \texttt{glmnet} standardises by default,
but it is worth knowing what the coefficients refer to when you plot
them.

\subsection{Setup}\label{setup}

\begin{Shaded}
\begin{Highlighting}[]
\FunctionTok{library}\NormalTok{(tidyverse)}
\FunctionTok{library}\NormalTok{(glmnet)}
\FunctionTok{set.seed}\NormalTok{(}\DecValTok{42}\NormalTok{)}
\FunctionTok{theme\_set}\NormalTok{(}\FunctionTok{theme\_minimal}\NormalTok{(}\AttributeTok{base\_size =} \DecValTok{12}\NormalTok{))}
\end{Highlighting}
\end{Shaded}

\subsection{1. Goal}\label{goal}

Compare ridge, lasso, and elastic net on a simulated high-dimensional
regression problem and pick a final model.

\subsection{2. Approach}\label{approach}

\texttt{n\ =\ 100}, \texttt{p\ =\ 200}, block-correlated design with a
handful of true signals, to exercise both the sparsity of lasso and the
group-handling of elastic net.

\begin{Shaded}
\begin{Highlighting}[]
\NormalTok{Z }\OtherTok{\textless{}{-}} \FunctionTok{matrix}\NormalTok{(}\FunctionTok{rnorm}\NormalTok{(n }\SpecialCharTok{*}\NormalTok{ (p }\SpecialCharTok{/} \DecValTok{10}\NormalTok{)), n, p }\SpecialCharTok{/} \DecValTok{10}\NormalTok{)}
\NormalTok{X }\OtherTok{\textless{}{-}}\NormalTok{ Z[, }\FunctionTok{rep}\NormalTok{(}\FunctionTok{seq\_len}\NormalTok{(p }\SpecialCharTok{/} \DecValTok{10}\NormalTok{), }\AttributeTok{each =} \DecValTok{10}\NormalTok{)] }\SpecialCharTok{+} \FloatTok{0.2} \SpecialCharTok{*} \FunctionTok{matrix}\NormalTok{(}\FunctionTok{rnorm}\NormalTok{(n }\SpecialCharTok{*}\NormalTok{ p), n, p)}
\NormalTok{beta }\OtherTok{\textless{}{-}} \FunctionTok{c}\NormalTok{(}\FunctionTok{rep}\NormalTok{(}\DecValTok{2}\NormalTok{, }\DecValTok{5}\NormalTok{), }\FunctionTok{rep}\NormalTok{(}\SpecialCharTok{{-}}\FloatTok{1.5}\NormalTok{, }\DecValTok{5}\NormalTok{), }\FunctionTok{rep}\NormalTok{(}\DecValTok{0}\NormalTok{, p }\SpecialCharTok{{-}} \DecValTok{10}\NormalTok{))}
\NormalTok{y }\OtherTok{\textless{}{-}} \FunctionTok{as.numeric}\NormalTok{(X }\SpecialCharTok{\%*\%}\NormalTok{ beta }\SpecialCharTok{+} \FunctionTok{rnorm}\NormalTok{(n))}

\FunctionTok{tibble}\NormalTok{(}\AttributeTok{i =} \FunctionTok{seq\_along}\NormalTok{(beta), }\AttributeTok{beta =}\NormalTok{ beta) }\SpecialCharTok{|\textgreater{}}
  \FunctionTok{ggplot}\NormalTok{(}\FunctionTok{aes}\NormalTok{(i, beta)) }\SpecialCharTok{+} \FunctionTok{geom\_segment}\NormalTok{(}\FunctionTok{aes}\NormalTok{(}\AttributeTok{xend =}\NormalTok{ i, }\AttributeTok{yend =} \DecValTok{0}\NormalTok{)) }\SpecialCharTok{+}
  \FunctionTok{labs}\NormalTok{(}\AttributeTok{x =} \StringTok{"feature index"}\NormalTok{, }\AttributeTok{y =} \StringTok{"true coefficient"}\NormalTok{)}
\end{Highlighting}
\end{Shaded}

\subsection{3. Execution}\label{execution}

\begin{Shaded}
\begin{Highlighting}[]
\NormalTok{fit\_ridge }\OtherTok{\textless{}{-}} \FunctionTok{glmnet}\NormalTok{(X, y, }\AttributeTok{alpha =} \DecValTok{0}\NormalTok{)}
\NormalTok{fit\_lasso }\OtherTok{\textless{}{-}} \FunctionTok{glmnet}\NormalTok{(X, y, }\AttributeTok{alpha =} \DecValTok{1}\NormalTok{)}
\NormalTok{fit\_enet  }\OtherTok{\textless{}{-}} \FunctionTok{glmnet}\NormalTok{(X, y, }\AttributeTok{alpha =} \FloatTok{0.5}\NormalTok{)}
\FunctionTok{par}\NormalTok{(}\AttributeTok{mfrow =} \FunctionTok{c}\NormalTok{(}\DecValTok{1}\NormalTok{, }\DecValTok{3}\NormalTok{), }\AttributeTok{mar =} \FunctionTok{c}\NormalTok{(}\DecValTok{4}\NormalTok{, }\DecValTok{4}\NormalTok{, }\DecValTok{2}\NormalTok{, }\DecValTok{1}\NormalTok{))}
\FunctionTok{plot}\NormalTok{(fit\_ridge, }\AttributeTok{xvar =} \StringTok{"lambda"}\NormalTok{); }\FunctionTok{title}\NormalTok{(}\StringTok{"ridge"}\NormalTok{)}
\FunctionTok{plot}\NormalTok{(fit\_lasso, }\AttributeTok{xvar =} \StringTok{"lambda"}\NormalTok{); }\FunctionTok{title}\NormalTok{(}\StringTok{"lasso"}\NormalTok{)}
\FunctionTok{plot}\NormalTok{(fit\_enet,  }\AttributeTok{xvar =} \StringTok{"lambda"}\NormalTok{); }\FunctionTok{title}\NormalTok{(}\StringTok{"elastic net"}\NormalTok{)}
\end{Highlighting}
\end{Shaded}

\subsection{4. Check}\label{check}

Pick \texttt{lambda} by CV for each; compare minimum CV MSE.

\begin{Shaded}
\begin{Highlighting}[]
\NormalTok{cv\_lasso }\OtherTok{\textless{}{-}} \FunctionTok{cv.glmnet}\NormalTok{(X, y, }\AttributeTok{alpha =} \DecValTok{1}\NormalTok{,   }\AttributeTok{nfolds =} \DecValTok{5}\NormalTok{)}
\NormalTok{cv\_enet  }\OtherTok{\textless{}{-}} \FunctionTok{cv.glmnet}\NormalTok{(X, y, }\AttributeTok{alpha =} \FloatTok{0.5}\NormalTok{, }\AttributeTok{nfolds =} \DecValTok{5}\NormalTok{)}
\FunctionTok{tibble}\NormalTok{(}
  \AttributeTok{model =} \FunctionTok{c}\NormalTok{(}\StringTok{"ridge"}\NormalTok{, }\StringTok{"lasso"}\NormalTok{, }\StringTok{"enet"}\NormalTok{),}
  \AttributeTok{min\_cvmse =} \FunctionTok{c}\NormalTok{(}\FunctionTok{min}\NormalTok{(cv\_ridge}\SpecialCharTok{$}\NormalTok{cvm), }\FunctionTok{min}\NormalTok{(cv\_lasso}\SpecialCharTok{$}\NormalTok{cvm), }\FunctionTok{min}\NormalTok{(cv\_enet}\SpecialCharTok{$}\NormalTok{cvm)),}
  \AttributeTok{nonzero   =} \FunctionTok{c}\NormalTok{(}\FunctionTok{sum}\NormalTok{(}\FunctionTok{coef}\NormalTok{(cv\_ridge, }\AttributeTok{s =} \StringTok{"lambda.min"}\NormalTok{) }\SpecialCharTok{!=} \DecValTok{0}\NormalTok{),}
                \FunctionTok{sum}\NormalTok{(}\FunctionTok{coef}\NormalTok{(cv\_lasso, }\AttributeTok{s =} \StringTok{"lambda.min"}\NormalTok{) }\SpecialCharTok{!=} \DecValTok{0}\NormalTok{),}
                \FunctionTok{sum}\NormalTok{(}\FunctionTok{coef}\NormalTok{(cv\_enet,  }\AttributeTok{s =} \StringTok{"lambda.min"}\NormalTok{) }\SpecialCharTok{!=} \DecValTok{0}\NormalTok{))}
\NormalTok{)}
\end{Highlighting}
\end{Shaded}

\subsection{5. Report}\label{report}

\begin{quote}
On a simulated \texttt{n\ =\ 100}, \texttt{p\ =\ 200} regression with 10
true signals, elastic net (α = 0.5) achieved the lowest 5-fold CV MSE
(\texttt{round(min(cv\_enet\$cvm),\ 2)}) and retained
\texttt{sum(coef(cv\_enet,\ s\ =\ "lambda.min")\ !=\ 0)\ -\ 1} features.
Lasso was comparable but dropped one member of each correlated signal
pair; ridge was smoothest but produced no sparse shortlist.
\end{quote}

The choice is problem-dependent: if the next step is validation of a
candidate biomarker panel, lasso or elastic net is preferable for the
interpretable sparsity; if the next step is prediction only, ridge is
often more stable.

\subsection{Common pitfalls}\label{common-pitfalls}

\begin{itemize}
\tightlist
\item
  Interpreting lasso coefficients as unbiased effect estimates.
\item
  Forgetting that correlated features share shrinkage: lasso may pick
  any one of them.
\item
  Reporting CV-minimum error as generalisation error (see the
  cross-validation lab).
\end{itemize}

\subsection{Further reading}\label{further-reading}

\begin{itemize}
\tightlist
\item
  Friedman J, Hastie T, Tibshirani R (2010), \emph{Regularization paths
  for generalized linear models via coordinate descent}.
\item
  Zou H, Hastie T (2005), \emph{Regularization and variable selection
  via the elastic net}.
\end{itemize}

\subsection{Session info}\label{session-info}

\begin{Shaded}
\begin{Highlighting}[]

\end{Highlighting}
\end{Shaded}

\subsection{See also --- chapter index}\label{see-also-chapter-index}

\begin{itemize}
\tightlist
\item
  \href{../index.Rmd}{Methods in this chapter}
\item
  \href{../../../ch00-find-your-method.Rmd}{Find your method (Chapter
  0)}
\end{itemize}

\end{document}
